\documentclass[12pt, a4paper]{article}

% ---- Ajustes de Idioma e Fontes Modernas -------------------
\usepackage[utf8]{inputenc}
\usepackage[T1]{fontenc}
\usepackage[brazil]{babel}
\usepackage{mathpazo}
\usepackage[scaled=0.95]{helvet}
\usepackage{courier}

% ---- Pacotes de Layout e Design ----------------------------
\usepackage{geometry}
\geometry{a4paper, top=3cm, bottom=2cm, left=3cm, right=2cm}
\usepackage{setspace}
\onehalfspacing
\usepackage{indentfirst}
\usepackage{microtype}

% ---- Cores e Hiperlinks ------------------------------------
\usepackage{xcolor}
\definecolor{brandblue}{RGB}{44, 62, 80}
\definecolor{softgray}{RGB}{245, 245, 245}
\usepackage[colorlinks=true, allcolors=brandblue]{hyperref}

% ---- Elementos Matemáticos e Gráficos ----------------------
\usepackage{amsmath, amssymb, amsfonts}
\usepackage{graphicx}
\usepackage{booktabs}
\usepackage{caption}
\usepackage{subcaption}
\usepackage{float}

% ---- Customização de Títulos -------------------------------
\usepackage{titlesec}
\titleformat{\section}{\color{brandblue}\normalfont\Large\bfseries}{\thesection}{1em}{}
\titleformat{\subsection}{\color{brandblue}\normalfont\large\bfseries}{\thesubsection}{1em}{}

% =============================================================
\begin{document}
% =============================================================

% ---- CAPA ---------------------------------------------------
\begin{titlepage}
    \begin{flushleft}
        \Large \textbf{Instituto Federal do Ceará} \\
        \large Mestrado em Ciência da Computação \\
        \small Disciplina: Reconhecimento de Padrões
    \end{flushleft}
    \vspace{4cm}
    \begin{center}
        \hrule\vspace{0.5cm}
        {\huge \textbf{\color{brandblue} Classificadores Bayesianos}}\\
        \vspace{0.2cm}
        {\Large Trabalho 03: Naive Bayes}
        \vspace{0.5cm}\hrule
    \end{center}
    \vspace{3cm}
    \begin{flushright}
        \large \textbf{Raquel Maciel Coelho de Sousa}\\
        \small Professor: Dr. Ajalmar Rocha Neto
    \end{flushright}
    \vfill
    \begin{center}\small Fortaleza, 11 de maio de 2026 \end{center}
\end{titlepage}

\newpage
{\hypersetup{linkcolor=black}\tableofcontents}
\newpage

% =============================================================
\section{Introdução}
% =============================================================

Este relatório descreve a implementação e avaliação do
\textbf{Classificador Naive Bayes Gaussiano}, desenvolvido como
complementação ao Trabalho~02 da disciplina de Reconhecimento de
Padrões. Diferentemente do Bayesiano Multivariado, o Naive Bayes
assume \textbf{independência condicional} entre os atributos dado
a classe, simplificando o cálculo da verossimilhança.

O modelo foi treinado e testado sobre os mesmos cinco conjuntos de
dados: Flor Íris, Coluna Vertebral, Breast Cancer Wisconsin,
Dermatology e Artificial~I, e seus resultados foram comparados com
o Bayesiano Multivariado, KNN e DMC.

% =============================================================
\section{Fundamentação Teórica}
% =============================================================

\subsection{Teorema de Bayes}

O ponto de partida é o mesmo dos demais classificadores:

\begin{equation}
    P(C_i \mid \mathbf{x}) =
    \frac{p(\mathbf{x} \mid C_i)\,P(C_i)}{p(\mathbf{x})}
\end{equation}

\noindent onde $P(C_i)$ é a probabilidade \textit{a priori},
$p(\mathbf{x} \mid C_i)$ é a verossimilhança e $p(\mathbf{x})$
é a evidência normalizadora.

\subsection{Pressuposto de Independência}

A diferença fundamental do Naive Bayes está na modelagem da
verossimilhança. Assumindo que os atributos são
\textbf{condicionalmente independentes} dado a classe $C_i$:

\begin{equation}
    p(\mathbf{x} \mid C_i) =
    \prod_{j=1}^{d} p(x_j \mid C_i)
    \label{eq:independencia}
\end{equation}

Ou seja, a verossimilhança conjunta é o \textbf{produto} das
verossimilhanças individuais de cada atributo. Isso elimina a
necessidade de estimar a matriz de covariância — cada atributo é
modelado por uma gaussiana univariada independente.

\subsection{Gaussiana Univariada por Atributo}

Cada fator $p(x_j \mid C_i)$ é modelado por uma gaussiana 1D:

\begin{equation}
    p(x_j \mid C_i) =
    \frac{1}{\sqrt{2\pi\sigma_{ij}^2}}
    \exp\!\left\{
        -\frac{(x_j - \mu_{ij})^2}{2\sigma_{ij}^2}
    \right\}
\end{equation}

\noindent onde $\mu_{ij}$ e $\sigma_{ij}^2$ são a média e a variância
do atributo $j$ na classe $C_i$, estimadas diretamente dos dados de
treino. A posteriori final combina todos os atributos:

\begin{equation}
    P(C_i \mid \mathbf{x}) \propto
    P(C_i) \prod_{j=1}^{d}
    \frac{1}{\sqrt{2\pi\sigma_{ij}^2}}
    \exp\!\left\{
        -\frac{(x_j - \mu_{ij})^2}{2\sigma_{ij}^2}
    \right\}
\end{equation}

\subsection{Comparação com o Bayesiano Multivariado}

A tabela abaixo resume as diferenças estruturais entre os dois modelos:

\begin{table}[H]
    \centering
    \caption{Naive Bayes vs.\ Bayesiano Multivariado.}
    \begin{tabular}{@{}lcc@{}}
        \toprule
        \textbf{Aspecto} & \textbf{Naive Bayes} & \textbf{Multivariado} \\
        \midrule
        Verossimilhança   & Produto de gaussianas 1D & Gaussiana $d$-dimensional \\
        Parâmetros por classe & $2d$ (pois $\mu_j$, $\sigma^2_j$) & $\approx\frac{d^2}{2}$ (pois $\boldsymbol{\mu}$, $\Sigma$) \\
        Correlações       & Ignoradas & Capturadas em $\Sigma$ \\
        Singularidade     & Não ocorre & Possível (regularização) \\
        Custo computacional & Baixo & Alto (inversão de $\Sigma$) \\
        \bottomrule
    \end{tabular}
\end{table}

\subsection{Classificadores de Referência}

\textbf{DMC} — classifica pelo centroide mais próximo em distância
euclidiana: $\hat{C} = \arg\min_{C_i}\|\mathbf{x}-\boldsymbol{c}_i\|_2$.

\textbf{KNN} ($k=5$) — voto majoritário entre os $k$ vizinhos mais
próximos no conjunto de treino.

\subsection{Escolha do Par de Atributos para Visualização}

O par foi selecionado pelo critério de Fisher:

\begin{equation}
    \text{score}(a_1,a_2) =
    \text{Fisher}(a_1)+\text{Fisher}(a_2)
\end{equation}

% =============================================================
\section{Conjuntos de Dados}
% =============================================================

\subsection{Flor Íris}
150 amostras, 4 atributos contínuos, 3 classes balanceadas.
Dataset limpo, sem valores faltantes.

\subsection{Coluna Vertebral}
310 amostras, 6 atributos biomecânicos contínuos, 3 classes
ortopédicas. Completamente numérico.

\subsection{Breast Cancer Wisconsin}
569 amostras, 30 atributos numéricos, 2 classes. Versão
\textit{Diagnostic} (UCI id=17).

\subsection{Dermatology}
366 amostras, 34 atributos, 6 classes. Valores faltantes imputados
pela média.

\subsection{Artificial I}
Dataset sintético gerado com 3 classes gaussianas 2D
($n=50$ por classe, semente $42$):
\begin{itemize}
    \item \textbf{Círculo:} $\boldsymbol{\mu}=(1.0,\,4.0)$, $\sigma=0.5$
    \item \textbf{Triângulo:} $\boldsymbol{\mu}=(4.0,\,4.0)$, $\sigma=0.5$
    \item \textbf{Estrela:} $\boldsymbol{\mu}=(2.0,\,2.0)$, $\sigma=0.5$
\end{itemize}

% =============================================================
\section{Metodologia}
% =============================================================

Foram realizadas \textbf{25 realizações} independentes por dataset,
com divisão 80\% treino e 20\% teste. A matriz de confusão foi
extraída da realização com acurácia \textbf{mais próxima da média}.
As superfícies de decisão foram geradas pelo Naive Bayes sobre o
par de atributos selecionado pelo critério de Fisher.

% =============================================================
\section{Resultados e Discussão}
% ============================================================

\subsection{Acurácia Média e Desvio Padrão}

\begin{table}[H]
    \centering
    \caption{Resultados experimentais --- 25 realizações, divisão 80/20.}
    \label{tab:results}
    \begin{tabular}{@{}lcccc@{}}
        \toprule
        \textbf{Base de Dados} & \textbf{Naive Bayes} & \textbf{Bayesiano} & \textbf{KNN ($k=5$)} & \textbf{DMC} \\
        \midrule
        Íris             & $94.8 \pm 4.3$ & $97.5 \pm 3.0$ & $96.3 \pm 2.9$ & $90.9 \pm 4.9$ \\
        Artificial I     & $99.5 \pm 1.2$ & $99.2 \pm 1.4$ & $99.2 \pm 1.7$ & $99.2 \pm 1.4$ \\
        Coluna Vertebral & $83.0 \pm 4.4$ & $83.9 \pm 4.5$ & $82.3 \pm 3.3$ & $75.7 \pm 5.5$ \\
        Breast Cancer    & $93.8 \pm 1.9$ & $88.4 \pm 3.0$ & $93.6 \pm 2.2$ & $87.9 \pm 2.7$ \\
        Dermatology      & $84.9 \pm 4.4$ & $88.6 \pm 3.8$ & $86.2 \pm 2.9$ & $50.4 \pm 6.2$ \\
        \bottomrule
    \end{tabular}
\end{table}

\subsection{Matrizes de Confusão}
Matriz de confusão de cada dataset com o classificador Naive Bayes. Buscando a melhor representatividade, foram utilizados os registros da realização com a taxa de acerto mais próxima da acurácia para montar cada matriz.

\begin{figure}[H]
    \centering
    \includegraphics[width=0.7\textwidth]{iris-matrix-confusao.png}
    \caption{Matriz de confusão --- Íris.}
    \label{fig:conf_multi}
\end{figure}

\begin{figure}[H]
    \centering
    \includegraphics[width=0.7\textwidth]{vertebral-matrix-confusao.png}
    \caption{Matriz de confusão --- Vertebral Column.}
    \label{fig:conf_multi}
\end{figure}

\begin{figure}[H]
    \centering
    \includegraphics[width=0.5\textwidth]{breast-matrix-confusao.png}
    \caption{Matriz de confusão --- Breast Cancer.}
    \label{fig:conf_multi}
\end{figure}

\begin{figure}[H]
    \centering
    \includegraphics[width=0.7\textwidth]{dermatology-matrix-confusao.png}
    \caption{Matriz de confusão --- Dermatology.}
    \label{fig:conf_multi}
\end{figure}

\begin{figure}[H]
    \centering
    \includegraphics[width=0.7\textwidth]{artificial-matrix-confusao.png}
    \caption{Matrix de confusão --- Artificial I.}
    \label{fig:conf_multi}
\end{figure}

\subsection{Superfícies de Decisão}
Superfícies de decisão do Naive Bayes com gaussianas univariadas sobrepostas. Treino~($\bullet$), teste~($\times$).

\begin{figure}[H]
\centering
\includegraphics[width=1\textwidth]{iris-superficie-decisao.png}
\caption{Superfície de decisão --- Íris}
\end{figure}

\begin{figure}[H]
\centering
\includegraphics[width=1\textwidth]{artificial-superficie-decisao.png}
\caption{Superfície de decisão --- Artificial I}
\end{figure}

\begin{figure}[H]
\centering
\includegraphics[width=1\textwidth]{vertebral-superficie-decisao.png}
\caption{Superfície de decisão --- Vertebral Column.}
\end{figure}

\subsection{Discussão}

O \textbf{Naive Bayes} demonstrou desempenho surpreendentemente
competitivo dada sua simplicidade. Seu destaque foi o
\textbf{Breast Cancer} ($93.8\%$), onde superou o Bayesiano
Multivariado ($88.4\%$). Isso ocorre porque, com 30 atributos e
$\sim$455 amostras de treino, a estimativa da matriz de covariância
$30\times30$ é ruidosa — o pressuposto de independência do Naive
Bayes evita esse problema ao custo de ignorar correlações.

No \textbf{Artificial~I} ($99.5\%$), o Naive Bayes apresentou o
melhor resultado absoluto: como as classes foram geradas com covariância
isotrópica (desvio padrão igual em ambas as direções), a ausência de
correlações entre atributos é exatamente satisfeita — o pressuposto
de independência é válido neste caso.

Na \textbf{Íris} ($94.8\%$), o Naive Bayes fica ligeiramente abaixo
do Bayesiano Multivariado ($97.5\%$) e do KNN ($96.3\%$), pois os
atributos da pétala são correlacionados entre classes e essa informação
é ignorada.

No \textbf{Dermatology} ($84.9\%$), o Bayesiano Multivariado
($88.6\%$) leva vantagem por capturar a estrutura conjunta das 6 classes.
O DMC colapsa para $50.4\%$, evidenciando que centroides únicos não
capturam a complexidade desse dataset.

% =============================================================
\section{Conclusão}
% =============================================================

O Naive Bayes mostrou-se uma escolha eficaz e eficiente, especialmente
em alta dimensionalidade (Breast Cancer) onde o Bayesiano Multivariado
sofre com estimativas ruidosas da matriz de covariância. O pressuposto
de independência, embora raramente verdadeiro, produz fronteiras de
decisão robustas na prática.

Comparado aos demais: supera o DMC em todos os datasets; equipara-se
ao KNN; perde para o Bayesiano Multivariado quando as correlações entre
atributos são informativas (Íris, Dermatology). Sua principal vantagem
é o custo computacional reduzido — estima $2d$ parâmetros por classe
contra $\approx d^2/2$ do Multivariado.

\newpage
\begin{thebibliography}{9}
    \bibitem{duda} DUDA, R. O.; HART, P. E.; STORK, D. G\@.
        \textbf{Pattern Classification}. 2.~ed. Wiley-Interscience, 2001.
    \bibitem{bishop} BISHOP, C. M\@.
        \textbf{Pattern Recognition and Machine Learning}. Springer, 2006.
    \bibitem{uci} DUA, D.; GRAFF, C\@. UCI Machine Learning Repository.
        Disponível em: \url{https://archive.ics.uci.edu/ml}. Acesso em: \today.
\end{thebibliography}

\end{document}